\subsection{A conditional local-repair reduction for terminal localization}

\begin{theorem}[C32: \status{Proved under stated hypotheses}; repair reduction]
\label{sa08:C32}
Let $k$ be a commutative ring. Consider the following two kinds of
augmented strict simplicial PD-DGA targets, with their coefficient maps:
\begin{enumerate}
\item $T=A\times_k B$, the pullback of the augmentations of two
admissible theories;
\item $T=H(f,g)=A\times_{B\times B}B^I$, for parallel maps
$f,g:A\to B$ of admissible theories, where $I=\Delta[1]$ and
the maps to $B\times B$ are $(f,g)$ and the two endpoints.
\end{enumerate}
For each such $T$, choose the partial model $P_T\to T$ given by
the linear-disk construction of Statement~\ref{sa08:C21}.
That construction applies because it uses only eventual horizontal
acyclicity of the target rows.

Assume, for every target of both kinds, that for each prime $p$ there
is a factorization
\[
 (P_T)_p\longrightarrow E_{T,p}\longrightarrow T_p
\]
with $E_{T,p}$ admissible, extending the specified map to $T_p$,
and whose first map is an actual isomorphism in every weight and
simplex in form degrees $q<d_{T,p}$. Assume $d_{T,p}\to\infty$:
for each fixed $d$ only finitely many primes have $d_{T,p}\leq d$.

Then every pair of admissible theories has an admissible common
source, and their category localized at all morphisms is terminal.
These assertions hold for arbitrary $k$ without requiring flat terms.
They also hold in the flat category when $k=\Z$ and the assumed
repairs have flat terms. The support condition $q\leq w$, if imposed
on all the data, is preserved.

The repair hypothesis for homotopy equalizers is separate from the
one for pairs. Neither is asserted to follow from the constructions
already proved above.
\end{theorem}
\begin{proof}
\par\medskip\noindent\textbf{The targets and their partial models.}
In positive weights, the rows of $A\times_k B$ are the direct sum
of the rows of $A$ and $B$. Thus they are eventually horizontally
acyclic. Its eventual vertical cohomology is the constant module
$S^w\oplus S^w$ in degree zero only.

For $H(f,g)$, endpoint restriction $B^I\to B\times B$ is
surjective in each sufficiently high-weight row, and its kernel is
horizontally acyclic. The resulting exact sequence with quotient
$A$ proves eventual horizontal acyclicity of $H(f,g)$. Vertically,
Statement~\ref{sa08:C16} identifies the both-endpoint kernel's
cohomology with $S^w$ in degree one. The connecting map on
coefficients is $f-g=0$, because both maps fix $S$. Hence, in every
fixed eventual range,
\[
 H^0(H(f,g))\cong\underline{S^w},\qquad
 H^1(H(f,g))\cong\underline{S^w},\qquad
 H^j(H(f,g))=0\quad(j>1).
\]
These are natural identifications of simplicial modules. The target
has an augmentation and coefficient map but is generally not
admissible.

The construction in Statement~\ref{sa08:C21} attaches free linear
vertical disks along normalized horizontal cycles. Maps into a target
extend because the corresponding target row is acyclic; no vertical
comparison of that target enters the lifting. It therefore constructs
$P_T\to T$ for either target above, with free terms and uniform
horizontal bounds. Its independent vertical disk calculation is
unchanged: through degree $d$ the coefficient comparison is exact
after inverting $\lfloor(d+1)/2\rfloor!$. In particular
$(P_T)_\Q$ is admissible.

\par\medskip\noindent\textbf{Arithmetic replacement with constant vertical cohomology.}
We need a slight extension of the path argument. Let $T$ have
horizontally eventually acyclic rows, and suppose that in each fixed
cochain range its vertical cohomology is, in sufficiently large
weights, a constant simplicial module in every degree. Positive
cohomology need not vanish. The relative prism
\[
 (\Delta[n]\times I,\,\Delta[n]\times\{0\})
\]
has zero relative cohomology with every constant coefficient module.
The finite-range slice hypercohomology comparison in
Statement~\ref{sa08:C16} uses, in the required total degrees,
\[
 E_2^{a,b}=H^a(\Delta[n]\times I,
                  \Delta[n]\times\{0\};\mathcal H^b(T))
       \Longrightarrow H^{a+b}(P_0(T)_n).
\]
Every relevant coefficient sheaf $\mathcal H^b(T)$ is constant,
so all these relative groups vanish. Thus $P_0(T)$ is eventually
vertically acyclic. The bound is uniform in $n$:
each relevant high-weight row is acyclic on every slice, and the
first-quadrant calculation uses only finitely many cochain rows and
comparison degrees. The same argument identifies the vertical
cohomology of $T^I$ with that of $T$.

Set $R=\prod_p\Z_{(p)}$ and $R_\Q=R\otimes_\Z\Q$, and use
scalar subscripts as in Statement~\ref{sa08:C30}. Form
\[
 J(T)=T_\Q\mathop{\times}_{T_{R\Q},e_0}
             (T_{R\Q})^I
           \mathop{\times}_{e_1,T_{R\Q}}T_R.
\]
The pure arithmetic sequence of that statement gives an exact
sequence in every row and every simplex,
\[
 0\longrightarrow T\xrightarrow{\eta_T}J(T)
       \longrightarrow P_0(T_{R\Q})\longrightarrow0.
\]
The last map subtracts the constant initial path. Its surjectivity
is precisely arithmetic exactness applied to the final value of a
based path. Thus $\eta_T$ is a horizontal weak equivalence in every
row. By the preceding relative-prism argument it is also an eventual
vertical quasi-isomorphism in every fixed range. In particular
$J(T)$ again has constant eventual vertical cohomology; it need not
be admissible.

The simplicial mapping path
\[
 \widetilde T=T\mathop{\times}_{J(T),e_0}J(T)^I
           \xrightarrow{e_1}J(T)
\]
is a trivial fibration in every simplicial-module row, because
$\eta_T$ is a horizontal weak equivalence. Its kernel is horizontally
acyclic. Its projection to $T$ has based-path kernel, so is an
eventual vertical quasi-isomorphism by the same constant-cohomology
argument. The two evaluations on $J(T)^I$ induce the same vertical
cohomology isomorphism by the prism comparison, and evaluation at zero
on $\widetilde T$ is $\eta_T$ composed with projection to $T$.
Consequently $\widetilde T\to J(T)$ is an eventual
vertical quasi-isomorphism as well. Since it is termwise surjective,
its kernel is eventually vertically acyclic.

\par\medskip\noindent\textbf{Glue and obtain actual maps.}
Apply the finite-idempotent construction of
Statement~\ref{sa08:C30} to the assumed $E_{T,p}$. It gives an
admissible theory $E_R$ and compatible maps
\[
 (P_T)_R\longrightarrow E_R\longrightarrow T_R.
\]
The one-edge arithmetic path pullback of $(P_T)_\Q$ and $E_R$
over $(E_R)_\Q$ is an admissible theory $E$, by the same proof as
C30: both its endpoint theories are admissible, and the coefficient
connecting map is the pure arithmetic difference map. The exact
endpoint diagrams give $E\to J(T)$. No admissibility of $T$ is
used in constructing $E$.

Now put
\[
 H=E\times_{J(T)}\widetilde T.
\]
Its projection to $E$ is termwise surjective with the kernel just
proved horizontally acyclic in all rows and vertically acyclic in
every eventual finite range. Hence $H$ is admissible and has an
actual map $H\to T$. For $T=A\times_k B$ this is the required
common source. For $T=H(f,g)$ it gives $h:H\to A$ and a strict
simplicial PD-DGA homotopy from $fh$ to $gh$.

In the localization, the endpoint maps $B^I\to B$ are equal:
both are inverses to the constant-path map $B\to B^I$.
Therefore $fh=gh$ there; since $h$ is inverted, every parallel
pair $f,g$ is equal after localization.

\par\medskip\noindent\textbf{All loops.}
Repeated common-source formation gives a source $L$ for the finitely
many objects $X_i$ in any zigzag loop. Choose maps $r_i:L\to X_i$.
For each edge $u:X_i\to X_j$, the maps $ur_i$ and $r_j$ are
parallel, so are equal in the localization. Thus the edge is
$r_jr_i^{-1}$ there, or its inverse when traversed backwards.
Their product around the loop is the identity. The category is
nonempty, since the standard theory exists. Its localization is
therefore connected with trivial vertex groups, hence terminal.

All constructions preserve the two gradings. If $k=\Z$ and the
terms and repairs are flat, the targets $T$, path constructions,
kernels and fiber products used here are subgroups of torsionfree
abelian groups and are flat. This gives precisely the flat version
stated, without asserting arbitrary-base flatness.
\end{proof}

Statement~\ref{sa08:C32} removes infinite arithmetic descent and
loop gluing as separate tasks once its local repairs are constructed.
It leaves the actual coefficient-preserving, low-degree-preserving
local extension repair problem. The nonzero $H^1$ of the homotopy
equalizer is explicit; ordinary repairs into admissible targets do
not automatically supply repairs into that target.
